% Vietnamese translation for OI Public Library
% Source-File: IOI中国国家候选队论文/2004/胡伟栋.ppt
\documentclass[11pt]{article}
\usepackage{oipl-vn}

\title{Giảm dư thừa và tối ưu hóa thuật toán\\{\large Ghi chú theo slide}}
\author{Hu Weidong}
\date{IOI 2004}

\begin{document}
\maketitle

\origtitle{Reducing redundancy and algorithm optimization}
\sourcepath{IOI 2004 / Hu Weidong.ppt}

\section{Định hướng}

Tệp PPT cũ chứa nhiều công thức và hình động, phần chữ trích được chỉ cho
thấy vài mốc như ``Balance''. Bản ghi này dựa trên bài viết PDF cùng chủ đề:
tối ưu thuật toán bằng cách nhận ra phần tính toán lặp lại hoặc phần trạng
thái không đóng góp thêm thông tin.

\section{Dạng dư thừa thường gặp}

Dư thừa có thể xuất hiện ở nhiều mức:
\begin{itemize}
  \item nhiều trạng thái khác nhau nhưng tương đương về ảnh hưởng tới tương
  lai;
  \item nhiều phép thử cùng kiểm tra một điều kiện đã biết;
  \item một cấu trúc dữ liệu lưu quá nhiều chi tiết so với truy vấn cần trả
  lời;
  \item thuật toán duyệt cả không gian trong khi chỉ một biên nhỏ có thể tạo
  đáp án tốt hơn.
\end{itemize}

\section{Cách giảm}

Bài viết PDF dùng các ví dụ về phân tích số nguyên và bài giá trị phần thưởng
lớn nhất. Trong cả hai, bước quan trọng là xác định phần nào của thông tin
đầu vào tạo ra các trường hợp thật sự khác nhau. Sau đó ta gom nhóm, dùng
đối xứng, dùng biên, hoặc thay thế phép duyệt bằng công thức đã tiền xử lý.

\section{Ghi nhớ}

Tối ưu không chỉ là viết nhanh hơn. Nó thường bắt đầu bằng câu hỏi: ``hai
tình huống này có cần phân biệt trong các bước sau không?'' Nếu câu trả lời
là không, ta nên hợp nhất chúng trong trạng thái hoặc trong cấu trúc dữ liệu.

\end{document}
